\documentclass{article}

% ---------------------------------------------------------------------
% Self-contained NeurIPS Datasets & Benchmarks-style preprint.
% Uses only base LaTeX packages so it compiles without the neurips.sty
% file. To submit to NeurIPS D&B, drop the body into neurips_2025.sty
% and remove the geometry/font block below.
% ---------------------------------------------------------------------
\usepackage[letterpaper,margin=1in]{geometry}
\usepackage[utf8]{inputenc}
\usepackage[T1]{fontenc}
\usepackage{times}
\usepackage{amsmath,amssymb}
\usepackage{booktabs}
\usepackage{graphicx}
\usepackage{xcolor}
\usepackage{url}
\usepackage[colorlinks=true,linkcolor=black,citecolor=blue,urlcolor=blue]{hyperref}
\usepackage{caption}
\usepackage{authblk}
\usepackage{natbib}

\title{A Calibration and Selective-Prediction Benchmark for
Perturbation-Effect Uncertainty: Distribution-Free Methods Hold,
Parametric Priors Miscover, and a Genetics Taxonomy Does Not Help}

\author{Anonymous Author(s)}
\affil{Hackathon submission (Datasets \& Benchmarks track)}
\date{}

\begin{document}
\maketitle

\begin{abstract}
We introduce a rigorous, donor-blocked calibration benchmark for
uncertainty quantification (UQ) of single-cell perturbation effects,
built on a CD4$^{+}$ T-cell Perturb-seq assay and an independent K562
Perturb-seq assay. The benchmark holds the point predictor fixed across
all UQ methods so that any difference in marginal coverage is
attributable to the calibration layer alone, and it evaluates methods on
a leave-one-donor-out (LODO) split of $58{,}558$ test rows spanning
$2{,}591$ perturbations and four donors. Our headline result is an
\emph{honest tie}: six distribution-free methods (jackknife$+$, Mondrian,
weighted and split conformal, naive-global quantile, and quantile
regression) are statistically indistinguishable in mean absolute
marginal coverage error, and the best conformal method does not separate
from the strongest non-conformal distribution-free baseline (paired
cluster-bootstrap $\Delta$ confidence interval includes zero). The
benchmark does, however, resolve a clear and significant contrast:
distribution-free calibration decisively outperforms a parametric
Gaussian prior, which over-covers by roughly $6.7$ percentage points at
the $80\%$ level under donor shift. We further contribute (i) a
cross-assay transfer protocol showing that frozen quantiles fail
($4.8/7.6/17.1\%$ empirical coverage at nominal $80/90/95\%$) but that
re-fitting the identical method on the target assay restores near-nominal
coverage across six calibration fractions and $500$ resplits; (ii) a
distribution-free risk-control certificate that is valid per-assay yet
honestly vacuous at $n{=}4$ exchangeable donors; and (iii) a
pre-registered negative showing that conditioning conformal prediction on
an external eQTL taxonomy does not sharpen intervals. All results were
hash-frozen and pre-registered before scoring, an independent audit
re-derived the benchmark surface from the raw $16.87$~GB data with zero
discrepancy, and we openly report the self-retraction of two earlier
surrogate ``wins.'' This is a benchmark and methods contribution; no
biological target is validated in scope.
\end{abstract}

% =====================================================================
\section{Introduction}
% =====================================================================

Machine-learning models for single-cell perturbation prediction are
proliferating, but a growing body of evidence suggests that headline
accuracy numbers in this field are frequently inflated. Recent audits
report that deep-learning perturbation-effect predictors do not yet
outperform simple linear baselines \citep{ahlmanneltze2025deep}, and
systematic benchmarks find that mean baselines are competitive with
foundation models on standard scores \citep{wu2024perturbench}. Much of
this fragility traces to a single methodological gap: point-prediction
metrics reward memorization of population structure and say nothing about
whether a model \emph{knows what it does not know}. A prediction of a
perturbation effect that is confidently wrong is more dangerous than one
that is honestly uncertain, yet the community lacks a rigorous, shared
protocol for measuring calibrated uncertainty under the distribution
shifts that matter for this data, namely shifts across biological donors
and across assays.

We address this gap with a calibration and selective-prediction
benchmark rather than a new predictor. Our motivating application is
CD4$^{+}$ T-cell Perturb-seq relevant to autoimmune biology, but we
stress at the outset that this is a measurement contribution: we ask
whether the uncertainty around a perturbation-effect estimate is
correct and useful, not whether any specific gene is a validated disease
target. The disease context supplies the data and the donor structure;
it supplies no biological claim.

The benchmark is designed around one controlling principle. Because
different UQ methods can differ both in their point predictions and in
their calibration, comparisons in the literature often conflate the two.
We eliminate that confound by fixing an identical base point predictor
(a histogram gradient-boosted regressor) across every method, so that
the only quantity that varies is the interval-construction layer. On top
of this fixed surface we evaluate nine UQ methods under a donor-blocked
leave-one-donor-out (LODO) split, which is the assumption-faithful
granularity for the donor shift that practitioners actually face.

Our contributions are as follows.
\begin{itemize}
\item \textbf{A powered, point-prediction-controlled head-to-head of nine
UQ methods} on $58{,}558$ LODO test rows across four donors, with paired
cluster-bootstrap significance testing over $2{,}591$ perturbations.
The result is an \emph{honest tie} among six distribution-free methods,
and a significant, reproducible loss for parametric priors.
\item \textbf{A cross-assay transfer protocol} demonstrating that frozen
calibration quantiles fail on an independent assay but that the
underlying method transfers once recalibrated, characterized across six
calibration fractions and $500$ resplits.
\item \textbf{A distribution-free risk-control analysis} (RCPS /
Learn-then-Test) that yields a valid non-vacuous certificate per assay
and honestly reports vacuity at $n{=}4$ donors, quantifying the price of
exchangeability rather than hiding it.
\item \textbf{A pre-registered negative result} showing that conditioning
conformal prediction on an external cis-eQTL taxonomy does not produce
sharper intervals at matched coverage.
\item \textbf{An integrity protocol}: hash-frozen pre-registration before
scoring, an independent audit that re-derives the benchmark surface from
the raw data, and an explicit self-retraction of two earlier surrogate
results.
\end{itemize}

% =====================================================================
\section{Benchmark Design}
% =====================================================================

\subsection{Data and target}

The primary benchmark surface is derived from a CD4$^{+}$ T-cell
Perturb-seq assay in primary human T cells
\citep{schmidt2022crispr}. For each perturbation we define the effect
target as $y = \log(1 + \lVert z \rVert_2)$, where $z$ is the vector of
differential-expression $z$-scores over a shared readout gene panel with
the on-target gene excluded, and the $z$-scores are Welch two-sample
statistics computed against negative-control cells. The independent
transfer assay is the Norman et al. K562 genome-scale genetic-interaction
Perturb-seq dataset \citep{norman2019exploring} (GEO accession
GSE133344), scored with an identical target definition over a shared
panel of $9{,}780$ genes ($95.1\%$ of the primary panel).

\subsection{The donor-blocked LODO surface}

The benchmark evaluates on a leave-one-donor-out split with four donor
folds. Each fold holds out all cells from one donor for testing and
calibrates on the remaining three, so that no donor appears in both
calibration and test. The pooled test set comprises $58{,}558$ rows
spanning $2{,}591$ unique perturbations, with per-fold test sizes of
$14{,}640$, $14{,}639$, $14{,}639$, and $14{,}640$ rows. This design
targets the shift that matters in practice: a model calibrated on some
donors must produce trustworthy uncertainty on a new donor.

\subsection{The point-prediction control}

The central design choice is that an identical base point predictor and
identical point predictions are held fixed across all nine UQ methods.
Consequently, any difference in marginal coverage or interval width
reflects the calibration layer alone and not a difference in the
underlying regression. We verified that the committed split-conformal
method reproduces its receipted empirical coverage exactly
($0.7932/0.8955/0.9486$ at $80/90/95\%$) on this surface, confirming that
the benchmark harness faithfully re-scores the committed pipeline rather
than a surrogate.

\subsection{Methods compared}

We evaluate nine UQ methods that fall into three families. The
\emph{distribution-free conformal} family comprises split conformal,
Mondrian conformal conditioned on culture condition
\citep{vovk2005algorithmic}, weighted conformal
\citep{tibshirani2019conformal}, and the jackknife$+$
\citep{barber2021jackknife}. The \emph{distribution-free adaptive} family
comprises quantile regression, conformalized quantile regression (CQR)
\citep{romano2019conformalized}, and a naive-global quantile baseline.
The \emph{parametric} family comprises a parametric Gaussian interval and
a naive Gaussian ensemble-variance interval. All methods target marginal
coverage at the $80$, $90$, and $95\%$ levels.

\subsection{Evaluation metrics and significance testing}

The primary metric is the mean absolute marginal coverage error across
the three levels, defined as the mean of $\lvert \hat{c}_\ell -
\ell\rvert$ for $\ell \in \{0.80, 0.90, 0.95\}$, where $\hat{c}_\ell$ is
the empirical coverage. We report interval width as an efficiency metric
and area under the risk-coverage curve (AURC) for selective prediction.
Significance is assessed with a paired cluster bootstrap ($2{,}000$
replicates) that resamples the $2{,}591$ perturbations as clusters, which
respects the dependence induced by shared perturbations across rows and
yields confidence intervals on pairwise differences between methods.

% =====================================================================
\section{Results}
% =====================================================================

\subsection{The core result: an honest tie among distribution-free methods}

Table~\ref{tab:main} reports mean absolute coverage error and mean
interval width at $90\%$ for all nine methods. Six distribution-free
methods form a tight cluster in coverage error, from jackknife$+$ at
$0.0018$ through Mondrian and weighted conformal at $0.0037$,
naive-global at $0.0040$, split conformal at $0.0042$, quantile
regression at $0.0049$, and CQR at $0.0054$. The paired cluster bootstrap
confirms that this cluster is statistically indistinguishable: the
comparison between the best committed conformal method (Mondrian) and the
strongest non-conformal distribution-free baseline (quantile regression)
gives a $\Delta$ of $-0.00074$ with a $95\%$ confidence interval of
$[-0.0034, +0.0023]$, which includes zero. Similarly, jackknife$+$ versus
split ($[-0.0028, +0.0025]$), Mondrian versus split
($[-0.0014, +0.0007]$), and CQR versus split ($[-0.0021, +0.0031]$) all
straddle zero.

This is the honest headline. On this surface, conformal prediction does
not beat every prior-art distribution-free method; quantile regression,
itself a distribution-free adaptive method, ties it. We report the tie
because it is the truth of the measurement, and because a tie among
principled distribution-free methods is itself a useful result for
practitioners choosing a calibration layer.

\begin{table}[t]
\centering
\caption{Head-to-head of nine UQ methods on the donor-blocked LODO
surface ($58{,}558$ test rows, four donor folds). Point predictions are
held identical across all methods. Coverage error is the mean absolute
marginal coverage error across the $80/90/95\%$ levels. Width is the mean
interval width at $90\%$. Lower is better for both. The six
distribution-free methods (top block) are statistically indistinguishable
in coverage error; the two parametric methods (bottom block) miscover.}
\label{tab:main}
\begin{tabular}{lccc}
\toprule
Method & Family & Coverage error $\downarrow$ & Width @90\% $\downarrow$ \\
\midrule
Jackknife$+$          & distribution-free & \textbf{0.0018} & 0.4605 \\
Mondrian conformal    & distribution-free & 0.0037 & 0.4581 \\
Weighted conformal    & distribution-free & 0.0037 & 0.4589 \\
Naive-global quantile & distribution-free & 0.0040 & 0.4581 \\
Split conformal       & distribution-free & 0.0042 & 0.4578 \\
Quantile regression   & distribution-free & 0.0049 & 0.4869 \\
CQR                   & distribution-free & 0.0054 & 0.4862 \\
\midrule
Parametric Gaussian   & parametric        & 0.0339 & 0.5328 \\
Ensemble-variance Gaussian & parametric   & 0.7887 & 0.0362 \\
\bottomrule
\end{tabular}
\end{table}

\subsection{The significant contrast: parametric priors miscover under shift}

The benchmark does resolve a clear and significant difference: the
distribution-free family decisively outperforms the parametric family.
The parametric Gaussian interval over-covers, reaching $86.67\%$
empirical coverage at the nominal $80\%$ level, an over-coverage of $6.7$
percentage points, with a mean absolute coverage error of $0.0339$,
nearly an order of magnitude worse than the distribution-free cluster.
The paired cluster bootstrap places the split-versus-parametric
difference at $-0.0287$ with a $95\%$ confidence interval of
$[-0.0376, -0.0152]$, which excludes zero; the Mondrian-versus-parametric
difference is $-0.0292$ with interval $[-0.0373, -0.0159]$, also
excluding zero. The naive Gaussian ensemble-variance interval collapses
entirely, achieving roughly $8\%$ coverage at the $90\%$ level (mean
absolute coverage error $0.7887$), because its variance estimate badly
underestimates the residual spread.

Two secondary observations sharpen the picture. First, adaptive-width
methods buy no sharpness here: at approximately equal coverage, the CQR
$90\%$ interval is significantly \emph{wider} than split conformal by
$+0.028$ log$_2$FC units (interval $[+0.0278, +0.0291]$), so CQR does not
Pareto-dominate the simpler incumbent, a consequence of the residual
distribution being close to homoscedastic in these leakage-safe features.
Second, on selective prediction all methods sit in a narrow AURC band of
roughly $0.112$ to $0.120$, far above the oracle lower bound of $0.048$
and close to the no-skill random line of $0.118$, indicating that
uncertainty-based trust ordering is only weakly informative on this
surface and does not separate the methods.

\subsection{Cross-assay transfer: frozen quantiles fail, recalibration fixes}

Calibration guarantees are exchangeability-bound, so a natural benchmark
question is whether quantiles calibrated on one assay transfer to
another. We applied the frozen CD4$^{+}$ conformal quantiles unchanged to
the independent K562 assay. Coverage collapsed to $4.8/7.6/17.1\%$ at
nominal $80/90/95\%$: a decisive failure. A diagnostic decomposition
shows the failure is both location and scale, since the K562 residual
spread is $2.6\times$ wider than the CD4$^{+}$ calibration spread, so even
oracle mean-recentering leaves coverage at only $0.41/0.49/0.68$. We
report this as a negative rather than tuning it away; it is the expected
behavior of split conformal under distribution shift.

Crucially, the \emph{method} transfers even though the specific quantiles
do not. Re-fitting the identical split-conformal procedure on the K562
assay's own calibration split restores near-nominal coverage on a
held-out K562 test split: a primary $50/50$ split gives $0.849/0.887/0.943$,
and a robustness sweep over $500$ random resplits gives mean coverage of
$0.755/0.864/0.927$, all three within the pre-registered $\pm0.05$ band of
nominal. A frozen-base sensitivity variant that recalibrates only the
quantiles (keeping the CD4$^{+}$ base predictor fixed) is even tighter at
$0.792/0.886/0.943$, showing that the assay-bound quantity is the quantile
and not the base predictor. Figure-level robustness holds at every
calibration fraction from $0.2$ to $0.7$: at each fraction the $95\%$
resplit distribution of coverage brackets the nominal level. This yields
a portable, cross-dataset calibration protocol with an honest failure
mode and its fix.

\subsection{Distribution-free risk control, honestly scoped}

We ask whether a finite-sample distribution-free certificate on
miscoverage can be issued using Risk-Controlling Prediction Sets (RCPS)
\citep{bates2021rcps} and Learn-then-Test \citep{angelopoulos2021learn}.
The certificate is exactly as strong as its exchangeability precondition.
Under donor-blocked LODO the assumption-faithful exchangeable unit is the
donor, of which there are only four. At $n{=}4$ donors the
distribution-free upper confidence bound on miscoverage is
$0.74/0.64/0.59$ (Hoeffding) against targets of $0.20/0.10/0.05$: the
certificate is honestly \emph{vacuous}. A row-level version that appears
tight ($n{=}58{,}558$) rests on row exchangeability across the donor
boundary, which is precisely the assumption the benchmark flags as
violated, so we report it only as a false-precision contrast, not a
guarantee.

In the per-assay recalibrated regime, however, exchangeability holds by
construction under a random calibration/test split of the $105$ K562
single-gene perturbations, and the certificate becomes valid and
non-vacuous. RCPS certifies a half-width $\hat{\lambda}=0.527$ at the
$90\%$ level (upper confidence bound $0.099 \le 0.10$) with realized
held-out coverage of $0.943$; certificates at $80/90/95\%$ all hold with
realized test coverage $0.887/0.943/0.943$. The contribution here is to
make the price of exchangeability explicit and quantified: risk control
is deployable per assay but cannot certify the donor-shift setting at
$n{=}4$, and we say so rather than overclaiming.

\subsection{A pre-registered negative: a genetics taxonomy does not help}

A plausible idea for sharpening perturbation intervals is to condition
the conformal calibration on an external, orthogonal genetics channel,
partitioning perturbations by a cis-eQTL tier and running Mondrian
conformal per tier. We pre-registered this hypothesis and tested it on
the $612$ genetics-annotated LODO test rows (tier sizes $360/168/84$).
The genetics-conditioned Mondrian method does not sharpen intervals
relative to plain split conformal; if anything it makes them
\emph{wider}. The paired-bootstrap difference in mean interval width
(genetics-Mondrian minus split) is $+0.020$ at $80\%$ (CI
$[+0.001,+0.034]$, excludes zero), $+0.022$ at $90\%$ (CI
$[-0.016,+0.057]$, includes zero), and $+0.219$ at $95\%$ (CI
$[+0.032,+0.438]$, excludes zero). At two of the three levels the
external genetics taxonomy is thus \emph{significantly worse} than plain
split conformal, and at no level is it better. Coverage was held (tier
coverages of $0.897/0.887/0.845$ at the $90\%$ level, all with confidence
intervals overlapping nominal). We report this as a first-class negative
result: it saves other groups from pursuing an intuitive but empirically
empty (indeed counterproductive) idea, which is exactly the kind of
result a benchmark paper should publish. We note the analysis is scoped
to $n{=}612$ annotated rows and is therefore underpowered, so the
confidence intervals are wide; we report it either way per
pre-registration.

% =====================================================================
\section{Reproducibility and Integrity}
% =====================================================================

Because this field is under active scrutiny for inflated metrics, we
built integrity guarantees into the benchmark protocol rather than
treating them as an afterthought. Every result was pre-registered and
hash-frozen before scoring: the primary benchmark carries a
pre-registration SHA-256 fingerprint, and the recalibration and
risk-control analyses each carry their own pre-registration hashes with
success bands declared before any number was computed. Receipts are
content-hashed, so any post-hoc alteration is detectable.

We commissioned an independent audit that re-derives the benchmark
surface directly from the raw $16.87$~GB single-cell data with fresh code
that does not reuse the pipeline's cache. For the audited donor fold, the
audit reproduces the effect-target multiset to within $10^{-6}$, with a
maximum sorted absolute difference of exactly $0.0$ and matching row
counts and feature columns, establishing that the cached benchmark
surface is a faithful materialization of the raw data and not a
fabrication. A second re-scoring audit confirms that the conformal
arithmetic on that surface is reproducible.

Finally, in the interest of full disclosure, we retract two earlier
surrogate results from this project's development history that did not
survive scrutiny. Reporting a self-retraction alongside a tie and a
negative result is uncomfortable, but it is the honest record of the
measurement, and it is the standard we believe perturbation UQ should be
held to.

% =====================================================================
\section{Limitations}
% =====================================================================

This benchmark has real limits that bound its claims. First, the
donor-blocked design has only \emph{four donors}, and the
exchangeable unit for the donor-shift guarantee is the donor. This is
adequate for measuring marginal calibration under LODO, which is our
purpose, but it is underpowered for any inferential claim that requires
more than a handful of exchangeable units, including the risk-control
certificate in the LODO regime, which is honestly vacuous at $n{=}4$. Any
extension that treats the donor as the inferential unit inherits this
ceiling.

Second, the genetics-conditioning experiment is a negative result scoped
to $612$ annotated rows and is underpowered; its confidence intervals are
wide, and we do not claim that no genetics-based conditioning could ever
help, only that this natural instantiation did not on this data.

Third, and most importantly, \emph{no biological outcome is validated in
scope}. This is a calibration and selective-prediction benchmark. The
CD4$^{+}$ T-cell and autoimmune context motivates the choice of data and
donor structure, but we make no claim that any perturbation is a
validated disease target. Calibrated uncertainty about a
perturbation-effect estimate is not evidence about biological efficacy.

Fourth, the transfer study uses a single external assay (K562), so the
recalibration protocol is demonstrated on one cross-assay pair rather
than a broad panel; generalization across many assays remains future
work.

% =====================================================================
\section{Conclusion}
% =====================================================================

We have presented a calibration and selective-prediction benchmark for
perturbation-effect uncertainty that leads with an honest tie:
distribution-free calibration methods are statistically
indistinguishable from one another on a fair, point-prediction-controlled
donor-blocked surface, and conformal prediction earns no measurable edge
over quantile regression. The benchmark nonetheless resolves a clear and
reproducible contrast, showing that parametric priors miscover under
donor shift while distribution-free methods do not. Alongside the core
head-to-head, we contribute a cross-assay transfer protocol with an
honest failure mode and its fix, a distribution-free risk-control
analysis that quantifies the price of exchangeability rather than hiding
it, and a pre-registered negative showing that an external genetics
taxonomy does not sharpen intervals. At a moment when the field is
alarmed about inflated perturbation metrics
\citep{ahlmanneltze2025deep, wu2024perturbench}, we hope a benchmark that
foregrounds its tie, its negative, its exchangeability limits, and its
self-retraction offers a useful and credible standard for measuring
calibrated uncertainty in single-cell perturbation prediction.

\paragraph{Availability.} All benchmark receipts are content-hashed and
deterministically reproducible from the committed pipeline, and the
external transfer analyses reproduce from cached group statistics for the
K562 assay.

\bibliographystyle{plainnat}
\bibliography{references}

\end{document}
